\section{Problem modeling}
\label{ch:problem_modeling}
This section introduces the signal model used throughout the thesis. Rather than attempting to describe the full physiological complexity of reinnervated EMG, the aim is to define a simplified mathematical framework that captures the main aspects relevant to source separation. The modeling approach is developed progressively. first, a general linear mixture model with multiple sources and observations is presented. After that, it is reduced to the two-source, two-channel case. Finally, it is specialized to the scenario of interest for the study, in which the target source is weaker and the donor-muscle contribution dominates one of the observed channels. For clarity, in the 2D cases, $\mathbf{s_1}$ denotes the \textbf{target flap-related source} and $\mathbf{s_2}$ the \textbf{donor-muscle source}. 
\clearpage

\subsection{General model}
The recorded EMG channels are modeled as mixtures of latent source signals. In a general setting, let $s_1(t), s_2(t),...,s_n(t)$ denote the underlying physiological sources and let $x_1(t),x_2(t),...,x_n(t)$ denote the observed channels. Each observation is expected to contain a weighted contribution of all sources, together with an additive noise term. This model can be expressed as a system of linear equations as shown in equation \ref{eq:inst_general_system}. Under the simplest instantaneous linear model, the observation vector can also be written as a linear transformation of the source vector plus noise, as shown in equation \ref{eq:inst_general_model}.
\begin{equation}
\left\{
\begin{aligned}
x_1(t) &= a_{11}s_1(t) + a_{12}s_2(t) + \cdots + a_{1n}s_n(t) + n_1(t) \\
x_2(t) &= a_{21}s_1(t) + a_{22}s_2(t) + \cdots + a_{2n}s_n(t) + n_2(t) \\
&\ \,\vdots \\
x_n(t) &= a_{n1}s_1(t) + a_{n2}s_2(t) + \cdots + a_{nn}s_n(t) + n_n(t)
\end{aligned}
\right.
\label{eq:inst_general_system}
\end{equation}
\begin{equation}
    x(t)=As(t)+n(t)
    \label{eq:inst_general_model}
\end{equation}

\noindent where $a_{ij}$ represents the contribution of source $j$ to observation $i$ and $A$ denotes the mixing matrix formed by the $a_{ij}$ coefficients.

This representation alone is not intended to fully describe the physical propagation of EMG through human tissue. Since the electrical activity generated by nearby muscles does not reach all electrodes under exactly the same conditions, each observation may contain contributions that are not only attenuated, but also slightly delayed with respect to one another. These delays may arise from differences in propagation path, electrode position, tissue properties and the spatial distribution of the underlying sources \cite{muceli_tutorial_2024}. As a result, the contribution of a given source to one channel does not necessarily appear at the same instant in another channel.

To account for this effect, the instantaneous mixture model can be extended by introducing delayed source terms. This leads to a convolutive mixture model, in which each observation depends on present and past values of the sources. However, for the purposes of this thesis, it is enough to only consider a single delay per channel and observation. This delay is implemented in equations \ref{eq:conv_general_system} and \ref{eq:conv_general_model}, which represent the extended convolutive version of \ref{eq:inst_general_system} and \ref{eq:inst_general_model}.

\begin{equation}
\left\{
\begin{aligned}
x_1(t) &= a_{11}s_1(t) + a_{12}s_2(t-\tau_{12}) + \cdots + a_{1n}s_n(t-\tau_{1n}) + n_1(t) \\
x_2(t) &= a_{21}s_1(t-\tau_{21}) + a_{22}s_2(t) + \cdots + a_{2n}s_n(t-\tau_{2n}) + n_2(t) \\
&\ \,\vdots \\
x_n(t) &= a_{n1}s_1(t-\tau_{n1}) + a_{n2}s_2(t-\tau_{n2}) + \cdots + a_{nn}s_n(t) + n_n(t)
\end{aligned}
\right.
\label{eq:conv_general_system}
\end{equation}

\begin{equation}
    x(t)=A_{\tau}s(t)+n(t)    
    \label{eq:conv_general_model}
\end{equation}

\noindent where $\tau_{ij}$ denotes the delay associated with  the contribution of source $j$ to travel to observation $i$. The notation  $A_{\tau}$ is used as a compact representation of the delayed mixing model.

While the previous equations describe the general $n$-source case, the remainder of this chapter focuses on a two-source, two-observation model. This reduced formulation is the simplest one in which source separation remains meaningful and it already captures the main features of the problem studied in this thesis. Equation \ref{eq:2d_conv_general_system} shows the general delayed model in the two-dimensional case.
\begin{equation}
\left\{
\begin{aligned}
x_1(t) &= a_{11}s_1(t) + a_{12}s_2(t-\tau_{12}) + n_1(t) \\
x_2(t) &= a_{21}s_1(t-\tau_{21}) + a_{22}s_2(t) + n_2(t) 
\end{aligned}
\right.
\label{eq:2d_conv_general_system}
\end{equation}

\subsection{Symmetric attenuation}
\label{sec:symmetric_attenuation}

The model introduced in the previous section can be considered as a standard \textit{Blind Source Separation (BSS)} setting, since no prior information is assumed about the values of the mixing coefficients $a_{ij}$ or about the properties of the original sources. However, given that this thesis aims to study specific scenarios related to the problem of interest, the current model proposed can be further improved by including additional information. 

In invasive EMG (or in surface EMG of relatively isolated muscles), each recorded channel can be approximated as the sum of the intended target source and an interfering contribution from the other source. Although invasive recordings are usually more selective than surface EMG, they do not guarantee perfect source isolation, since the adjunct muscular activity can still interfere with the measured channel 
\cite{m_consensus_nodate,l_analysis_2020}. Each channel is assumed to be dominated by its corresponding source, so the direct coefficients are simplified to $a_{ii}=1$. If both muscles are equally dominant an their interference is balanced, the off-diagonal terms can be expressed by a parameter $\beta$, where $a_{ij}=\beta$ for $\forall i\neq j$. Equation \ref{eq:2d_conv_beta_system} expresses the new model after taking into account these two considerations.

\begin{equation}
\left\{
\begin{aligned}
x_1(t) &= s_1(t) + \beta s_2(t-\tau_{12}) + n_1(t) \\
x_2(t) &= \beta s_1(t-\tau_{21}) + s_2(t) + n_2(t) 
\end{aligned}
\right.
\label{eq:2d_conv_beta_system}
\end{equation}

\subsection{Weak-source regime}
\label{sec:weak_source_regime}

Unfortunately, in reinnvervated muscle configurations, the interaction between the two muscular contributions is not necessarily balanced. In some cases, one source may display a stronger influence on the observed mixtures that the other, sot he coupling between channels cannot always be assumed to be symmetric. In fact, in procedures like VDMT, the contribution of the auxiliary muscle over the target flap-related source ($a_{12}$) is generally greater than the reverse interaction ($a_{21}$), which can barely be noticed in the auxiliary muscle observation. Mathematically, this can be expressed as $\mathbf{a_{21}=\alpha}$ with $\mathbf{\alpha<<1}$  (for most experiments $\mathbf{\alpha \approx0.01}$). Equation \ref{eq:2d_beta_asymmetric_system} reflects the mentioned piece of information.

\begin{equation}
\left\{
\begin{aligned}
x_1(t) &= s_1(t) + \beta s_2(t-\tau) + n_1(t) \\
x_2(t) &= \alpha s_1(t-\tau) + s_2(t) + n_2(t) 
\end{aligned}
\right.
\label{eq:2d_beta_asymmetric_system}
\end{equation}
\\
With these assumptions, this last model can may also be viewed as regression problem, where one channel consists of a dominant source and scaled interfering contribution that be measured as a reference. However, the central perspective of this thesis remains that of a signal separation problem, since the objective is to discuss general scenarios, even those where there is no reliable reference. A regression formulation will  be revisited later in the thesis for comparison, together with its main advantages and limitations.

One final note about this model is that since the tissue propagation is generally assumed to attenuate rather than amplify cross-talk contributions, the $\beta$ parameter is restricted to the interval $[0,1]$ region. However, it can be expected that the variance of the auxiliary signal $\mathrm{Var}(s_2)$ is greater than the variance of the target flap-related source $\mathrm{Var}(s_1)$, which will favor the retrieve of the auxiliary muscle source.  

